\section{Equilibrium and Tariff Architecture}\label{sec:equilibrium}

\subsection{Continuation pricing on the certified both-served branch}

Fix an installed state $(l,h)$ in the candidate interval derived in Section~\ref{sec:model} and covered by strict $\mathcal R^+$. Weak participation at zero surplus allows the provider to set
\begin{equation}
 F=S_L(p)=\frac{(a_L-p)^2}{2},
 \label{eq:Fbind}
\end{equation}
while serving both classes. Type $L$ then receives zero continuation utility and participates by convention. Type $H$ retains rent
\begin{equation}
 R_H(p)=S_H(p)-S_L(p)=R_F-dp.
 \label{eq:Hrent}
\end{equation}
Under flat pricing, $p=0$ and type $H$ receives $R_F$.

The global continuation proof is given in Appendix~\ref{app:proofs}; it does not rely on an interior FOC alone. Conditional on both types being served, differentiation gives
\begin{equation}
 \frac{\partial \Pi_M^B}{\partial p}=dh-(l+h)(p-c),
 \qquad
 \frac{\partial^2 \Pi_M^B}{\partial p^2}=-(l+h)<0,
 \label{eq:metered-FOC-SOC}
\end{equation}
and the exact square-loss identity in the Appendix establishes the global both-served optimum
\begin{equation}
 p^*(h)=c+d\frac{h}{l+h}.
 \label{eq:pstar}
\end{equation}
Condition \eqref{eq:Rplus4} keeps this price strictly between zero and $a_L$ over the complete candidate interval. The remaining $\mathcal R^+$ conditions, together with the active-set comparison in the Appendix, make the both-served continuation the provider's global continuation optimum on that interval.

The gross profit advantage of optimized metering over optimized flat pricing on this branch is
\begin{equation}
 \Phi(l,h)
 =\frac{l+h}{2}\left(c+d\frac{h}{l+h}\right)^2.
 \label{eq:Phi}
\end{equation}
The provider meters at state $(l,h)$ when $\Phi(l,h)>\mu$, chooses flat when $\Phi(l,h)<\mu$, and is indifferent at equality. Moreover,
\begin{equation}
 \frac{\partial\Phi(l,h)}{\partial h}
 =\frac{[c(l+h)+dh][c(l+h)+dh+2dl]}{2(l+h)^2}>0.
 \label{eq:Phi-derivative}
\end{equation}
Thus a larger installed high-use mass makes metering more attractive to the provider.

\subsection{Architecture-sensitive integration}

Because the low-use participation constraint binds under either architecture on the certified branch, low-use integration depends only on the current non-expropriable benefit:
\begin{equation}
 l=\frac{b_L}{K_L}.
 \label{eq:lstate}
\end{equation}
For the high-use class, future continuation rent depends on anticipated architecture.

Under a flat expectation,
\begin{equation}
 h_F=\frac{b_H+R_F}{K_H}.
 \label{eq:hF}
\end{equation}
If instead users expect metering, the rational-expectations high-use mass $h_M$ solves
\begin{equation}
 K_H h=b_H+R_F-d\left(c+d\frac{h}{l+h}\right).
 \label{eq:hMfixedpoint}
\end{equation}
The state $h_0=b_H/K_H$ is the integration level when no future continuation rent is expected.

\begin{proposition}[Integration ordering]\label{prop:T1}
For every parameter vector in strict $\mathcal R^+$, equation \eqref{eq:hMfixedpoint} has a unique solution and
\begin{equation}
 h_0<h_M<h_F.
 \label{eq:T1}
\end{equation}
\end{proposition}

The first inequality is substantive after the continuation repair: metering leaves the high-use class positive continuation rent on the certified both-served branch, while the second follows because metering reduces that rent relative to flat pricing. The proof in Appendix~\ref{app:proofs} uses the strict fixed-point crossing implied by $\mathcal R^+$.

\subsection{Threshold separation, boundary equilibria, and the pure-regime gap}

Define the expectation-contingent architecture thresholds
\begin{equation}
 \mu_M\equiv\Phi(l,h_M),
 \qquad
 \mu_F\equiv\Phi(l,h_F).
 \label{eq:mu-thresholds}
\end{equation}

\begin{proposition}[Threshold separation]\label{prop:T2}
For every parameter vector in strict $\mathcal R^+$,
\begin{equation}
 \mu_M<\mu_F.
 \label{eq:T2}
\end{equation}
\end{proposition}

This is the central state-feedback comparison: anticipated architecture changes the installed state, and the provider's gross metering gain is strictly ordered in that state. The abstract implication ``ordered states plus an increasing gain map imply ordered thresholds'' is only an organizing lemma; the economic content of Proposition~\ref{prop:T2} comes from the primitive derivation of both orderings in the quadratic baseline.

Under a pure metered expectation the provider evaluates architecture at $h_M$. Because provider indifference is allowed, the pure metered outcome is self-consistent if and only if
\begin{equation}
 \mu\leq\mu_M.
 \label{eq:metered-pure-condition}
\end{equation}
Under a pure flat expectation it evaluates architecture at $h_F$, so the pure flat outcome is self-consistent if and only if
\begin{equation}
 \mu\geq\mu_F.
 \label{eq:flat-pure-condition}
\end{equation}
At $\mu=\mu_M$ the aggregate pure equilibrium has $h=h_M$ and metering probability $\rho=1$; at $\mu=\mu_F$ it has $h=h_F$ and $\rho=0$. Provider indifference at those states does not make every mixing probability aggregate-consistent: any $\rho<1$ at $\mu_M$ raises high-use integration above $h_M$, making metering strictly optimal, while any $\rho>0$ at $\mu_F$ lowers high-use integration below $h_F$, making flat pricing strictly optimal. Appendix~\ref{app:proofs} gives the formal best-response and aggregation argument.

\begin{proposition}[Pure-regime gap]\label{prop:T3}
For every parameter vector in strict $\mathcal R^+$ and every $\mu$ satisfying
\begin{equation}
 \mu_M<\mu<\mu_F,
 \label{eq:strict-gap}
\end{equation}
neither pure flat nor pure metered expectation is self-consistent. On the certified branch there is a unique aggregate mixed resolution $(h^*,\rho^*)$ with $h_M<h^*<h_F$ and $0<\rho^*<1$.
\end{proposition}

The mixed state is the unique solution to
\begin{equation}
 \Phi(l,h^*)=\mu.
 \label{eq:hstar}
\end{equation}
At $h^*$ the provider is indifferent between the two \emph{actual globally optimal tariffs}: the flat both-served tariff $(F,p)=(a_L^2/2,0)$ and the metered both-served tariff $(F,p)=(S_L(p^*(h^*)),p^*(h^*))$. The provider randomizes between these tariffs, and $\rho$ denotes the probability that the metered tariff is realized. Users integrate before that randomization and therefore compare expected realization-contingent continuation utilities; they do not choose future quantities at an average price.

If the provider meters with probability $\rho$, expected high-use continuation rent is $R_F-\rho d p^*(h^*)$, so aggregate integration requires
\begin{equation}
 \rho^*=\frac{K_H(h_F-h^*)}{d\,p^*(h^*)},
 \qquad 0<\rho^*<1.
 \label{eq:rhostar}
\end{equation}
A strategy construction is
\begin{equation}
 x_L(k)=\mathbf 1\{k\leq b_L\},
 \qquad
 x_H(k)=\mathbf 1\{k\leq b_H+R_F-\rho^*d p^*(h^*)\}.
 \label{eq:integration-strategies}
\end{equation}
Under the uniform cost distributions these cutoffs aggregate to $l=b_L/K_L$ and $h=h^*$. A single atomless user's deviation does not change $(l,h)$ and therefore does not alter the provider's continuation problem. At off-path installed states the provider follows the global continuation values $V_F$ and $V_M$ in Section~\ref{sec:model}, mixing only at a tie. The uniqueness claim in Proposition~\ref{prop:T3} is intentionally aggregate: measure-zero cutoff behavior and off-path tie-breaking need not be unique, and no independent mixing by the continuum of users is required.

\begin{figure}[ht]
  \centering
  \input{../figures/threshold_phase.tex}
  \caption{Two self-consistency thresholds in the frozen exact example. Pure metering is self-consistent through the left endpoint $\mu_M$, pure flat pricing is self-consistent from the right endpoint $\mu_F$, and only the open middle interval is the mixed pure-regime gap. The numerical endpoints are illustrative; Propositions~\ref{prop:T1}--\ref{prop:T3} are quantified only on strict $\mathcal R^+$.}
  \label{fig:threshold-phase}
\end{figure}

\begin{table}[ht]
  \centering
  \input{../tables/baseline_example.tex}
  \caption{Generated exact objects for the permanent strict-$\mathcal R^+$ example. The table is produced by the repository generator and serves as a numerical illustration rather than a substitute for the theorem quantifiers.}
  \label{tab:baseline-example}
\end{table}

\subsection{Architecture-insensitive integration benchmark}

\begin{corollary}[Threshold collapse]\label{cor:collapse}
If integration is architecture-insensitive so that $h_M=h_F$, then
\begin{equation}
 \mu_M=\mu_F.
\end{equation}
The strict pure-regime gap disappears.
\end{corollary}

This benchmark isolates the mechanism. The gap is not generated merely by a two-part tariff or by the provider's inability to commit. It requires expected architecture to change the installed state that later changes the provider's own architecture incentive.
